\section{Topology of the protected-state bundle}
\label{sec:topology}

Local curvature statistics do not exhaust the information carried by a
degenerate state bundle.  The bundle can be locally close to a random matrix
field and still obey global quantization conditions imposed by the topology of
the moduli space.  This distinction is central to the geometric formulation:
local eigenvalue repulsion is a statement about a matrix at one point, whereas
a Chern number is an integral constraint on a family of matrices over a closed
cycle.  The two statements can coexist, but neither implies the other.

\subsection{The bundle and its gauge-covariant curvature}

Let $\mathcal{M}_{\mathrm{reg}}$ denote a connected region of coupling space
on which the protected rank is constant and the external gap remains open.
The BPS eigenspaces form a rank-$D$ complex vector bundle
\begin{equation}
 \mathcal{E}_{\mathrm{BPS}}\longrightarrow\mathcal{M}_{\mathrm{reg}},
 \qquad
 \mathcal{E}_{\mathrm{BPS}}|_{\lambda}=\operatorname{im}P(\lambda).
 \label{eq:topological-bps-bundle}
\end{equation}
Choose a local orthonormal frame $|a(\lambda)\rangle$ for the fiber.  The
Berry connection and curvature are
\begin{equation}
 (A_\mu)_{ab}=i\langle a|\partial_\mu b\rangle,
 \qquad
 F_{\mu\nu}=\partial_\mu A_\nu-\partial_\nu A_\mu-i[A_\mu,A_\nu].
 \label{eq:topological-connection}
\end{equation}
Under a frame change $|a\rangle\mapsto|b\rangle U_{ba}(\lambda)$, one has
$A\mapsto U^{-1}AU+iU^{-1}dU$ and $F\mapsto U^{-1}FU$.  Consequently,
$\operatorname{Tr}F$ and the eigenvalues of $F$ are frame independent,
although the individual matrix entries are not.  This is the non-Abelian
extension of the adiabatic phase and the geometric metric construction
\cite{berry1984,wilczekzee1984,provostvallee1980}.  Kato's projector
transport gives the corresponding formulation without choosing a frame
\cite{kato1950}.

The topology is meaningful only after the domain has been specified.  If a
submanifold of $\mathcal{M}$ contains a point where an additional BPS state
enters the protected sector, or where the external gap closes, the rank-$D$
bundle is not defined across that point.  One must excise the singular locus
and work on $\mathcal{M}_{\mathrm{reg}}$.  Treating the excised locus as an
ordinary point would confuse a change of bundle rank with a large but finite
curvature.

\subsection{Chern data as an integral constraint}

For a closed oriented two-cycle $\Sigma\subset\mathcal{M}_{\mathrm{reg}}$, the
first Chern number of the determinant line bundle is
\begin{equation}
 c_1(\mathcal{E}_{\mathrm{BPS}};\Sigma)
 =\frac{1}{2\pi}\int_{\Sigma}\operatorname{Tr}F\in\mathbb{Z},
 \label{eq:first-chern-bps}
\end{equation}
with the normalization in Eq.~\eqref{eq:topological-connection}.  The
many-body formulation of quantized Hall response established that the same
type of integral can be defined over a parameter torus even when the fiber is
an interacting many-body ground-state sector
\cite{niu1985}.  On a discretized parameter space, link variables
provide a gauge-invariant numerical implementation of the corresponding
integer \cite{fukui2005}.

Equation~\eqref{eq:first-chern-bps} is stronger than a statement about the
mean or variance of a local curvature entry.  Small perturbations of the
curvature that preserve the bundle and the cycle cannot change the integer.
Conversely, a change in the integer requires a singularity, a change of cycle,
or a loss of the rank-protection and gap assumptions.  This is why a random
matrix comparison must be made together with a topology audit: a genuinely
independent field of matrices would not, without additional constraints,
reproduce the quantized flux of a nontrivial bundle.

For a flat connection on a contractible patch, $F=0$ implies that the local
first Chern integral vanishes.  This local implication should not be promoted
to a statement that the entire moduli space is topologically trivial.  A flat
connection may have nontrivial holonomy on a non-simply-connected domain, and
the domain itself may acquire nontrivial cycles after singular strata are
removed.

\subsection{Topology in the \texorpdfstring{$\mathcal{N}=2$}{N=2} SYK moduli space}

The direct topological calculation relevant here is the $\mathcal{N}=2$ SYK
analysis of Ref.~\cite{chen2026berry}.  The paper identifies special loci in
the coupling space where additional BPS states appear, studies the regular
moduli space obtained after their removal, and integrates the Berry curvature
over candidate two-spheres.  It reports quantized Chern numbers and finds that
the generic-point curvature matrix can resemble a random matrix while its
integral remains constrained by Eq.~\eqref{eq:first-chern-bps}.

For the family of two-spheres and charge sectors analyzed there, the reported
first Chern number is
\begin{equation}
 c_1^{(r)}
 =\frac{1}{2\pi}\int_{\Sigma}\operatorname{Tr}F^{(r)}
 =\binom{N}{r+3}-\binom{N}{r-3}.
 \label{eq:syk-chern-formula}
\end{equation}
The formula is an empirical result of the stated moduli-space construction;
the source reports checks through $N=10$ and does not claim a systematic
classification of all cycles.  In particular, the even-$N$ sector $r=N/2$
has vanishing value in the formula and must not be included in a blanket
statement that every large-$N$ sector carries an exponentially large integer.

For the neighboring even sector $r=N/2-1$, Eq.~\eqref{eq:syk-chern-formula}
gives the large-$N$ estimate
\begin{equation}
 c_1^{(N/2-1)}
 \sim 24\sqrt{\frac{2}{\pi}}\,N^{-3/2}2^N,
 \qquad N\to\infty,
 \label{eq:syk-chern-asymptotic}
\end{equation}
as reported in Ref.~\cite{chen2026berry}.  The appropriate wording is thus
that the studied $\mathcal{N}=2$ SYK family contains exponentially large first
Chern numbers in specified large-$N$ sectors.  It is not correct to state,
without these qualifications, that all Chern numbers of $\mathcal{N}=2$ SYK are
exponential or that the asymptotic formula has been proved for every cycle.

The mechanism is transparent in the bundle language.  The singular strata
act as sources of Berry flux because the protected rank changes when they are
approached.  A two-sphere that avoids the singular strata can nevertheless
enclose them in the ambient coupling space, giving a nonzero integral of
$\operatorname{Tr}F$.  The local curvature away from the strata can therefore
show random-matrix-like eigenvalue statistics while the global flux remains an
integer.  This is a concrete form of structured geometric randomness: local
randomness and global topological memory are not mutually exclusive
\cite{chen2026berry}.

\subsection{Topology versus local random-matrix statistics}

There are three logically distinct tests for a protected bundle:
\begin{enumerate}
  \item a local test, such as eigenvalue repulsion or a distribution of
  curvature invariants at one point;
  \item a transport test, such as a Wilson loop or a non-Abelian holonomy; and
  \item a global test, such as a Chern number on a specified cycle.
\end{enumerate}
The Jacobi and random-compression ensembles provide useful local reference
models for the first test \cite{collins2005,zyczkowski2000}.  They do not
encode the topology of the parameter space.  The Chern integral supplies an
independent check that can reject a naive identification of the curvature
field with an unconstrained collection of independent random matrices.

This point also limits the interpretation of non-Abelian curvature in the
smooth sectors.  If the direct calculation finds $F=0$ for a particular
family, then the first Chern integral over any cycle contained in that flat
region is zero.  The result does not exclude nontrivial topology elsewhere in
the full moduli space, nor does it imply that a different BPS multiplet has a
flat connection.  The source-specific sector and deformation prescription
must accompany every such statement \cite{chen2026berry}.

Likewise, a random-matrix-like local spectrum in $\mathcal{N}=2$ SYK or
super-JT does not erase the Chern constraint.  A geometric response ansatz
that is intended to describe a full moduli space must preserve both properties:
the local distribution of $F$ and the integral quantization of its curvature.
The next section formulates this as a conditional prediction rather than as a
theorem about BPS black holes.
